% 数学ノート集・共通テンプレート v2（積分論2026年度第1部の版面を参考）
% A5判。章内で定義・定理・例・注意等の番号を共通化。
% コンパイル： platex -interaction=nonstopmode math_note_template_v2.tex (2回)
%              dvipdfmx math_note_template_v2.dvi
\documentclass[a5paper,10pt,onecolumn]{jsarticle}
\usepackage[dvipdfmx,top=22mm,bottom=20mm,inner=19mm,outer=18mm,headheight=13pt,headsep=6mm]{geometry}
\usepackage{amsmath,amssymb,amsthm,mathtools}
\usepackage{enumitem,fancyhdr,titlesec}
\usepackage[dvipdfmx,bookmarks=true,colorlinks=false]{hyperref}
\usepackage{pxjahyper}

% ----- 差し替える文書情報 -----
\newcommand{\NoteTitle}{凸最適化とKarush--Kuhn--Tucker条件}
\newcommand{\NoteSeries}{数学ノート集}

% ----- 文字・組版 -----
% 参考PDFの質素な明朝＋Computer Modernの印象を引き継ぎながら，
% A5判で窮屈にならないよう本文の行間を少し確保する．
\linespread{1.13}
\setlength{\parindent}{1zw}
\setlength{\parskip}{0pt}
\setlength{\abovedisplayskip}{6pt plus 2pt minus 2pt}
\setlength{\belowdisplayskip}{6pt plus 2pt minus 2pt}
\setlength{\abovedisplayshortskip}{3pt plus 2pt minus 1pt}
\setlength{\belowdisplayshortskip}{4pt plus 2pt minus 1pt}
\setlist[enumerate]{topsep=3pt,itemsep=1pt,parsep=0pt,leftmargin=2.4zw}
\titleformat{\section}{\normalfont\normalsize\bfseries}{\thesection}{0.8zw}{}
\titlespacing*{\section}{0pt}{2.0ex plus .4ex}{.9ex}
\titleformat{\subsection}{\normalfont\normalsize\bfseries}{\thesubsection}{0.7zw}{}
\titlespacing*{\subsection}{0pt}{1.5ex plus .3ex}{.6ex}
\raggedbottom

% ----- 柱とページ番号：ページ上部、罫線なし -----
\fancypagestyle{mathnote}{%
  \fancyhf{}%
  \fancyhead[L]{\footnotesize \NoteSeries\ （\NoteTitle）}%
  \fancyhead[R]{\footnotesize \thepage}%
  \renewcommand{\headrulewidth}{0pt}%
  \renewcommand{\footrulewidth}{0pt}%
}
\pagestyle{mathnote}

% ----- 節内通し番号の定義・定理群 -----
\newtheoremstyle{mathnotetheorem}
 {4.5pt}{4.5pt}{\normalfont}{}%
 {\bfseries}{}{0.65zw}%
 {\thmname{#1}\thmnumber{ #2}\thmnote{（#3）}}
\theoremstyle{mathnotetheorem}
\newtheorem{theorem}{定理}[section]
\newtheorem{definition}[theorem]{定義}
\newtheorem{lemma}[theorem]{補題}
\newtheorem{proposition}[theorem]{命題}
\newtheorem{corollary}[theorem]{系}
\newtheorem{example}[theorem]{例}
\newtheorem{remark}[theorem]{注意}
\renewcommand{\proofname}{証明}
\renewcommand{\qedsymbol}{$\square$}
\makeatletter
\renewenvironment{proof}[1][\proofname]{%
  \par\pushQED{\qed}\normalfont\topsep4\p@\@plus2\p@\relax
  \trivlist\item[\hskip\labelsep\bfseries #1]\ignorespaces
}{\popQED\endtrivlist\@endpefalse}
\makeatother

% ----- 使うノートで適宜拡張する記号 -----
\newcommand{\R}{\mathbb{R}}
\newcommand{\norm}[1]{\lVert #1\rVert}
\newcommand{\abs}[1]{\lvert #1\rvert}

\begin{document}
\section*{読了後にできること}
本稿を読むための前提は，凸関数，KKT条件，Slater条件．
本稿の具体的な到達課題は，不等式制約付き問題の最適解と乗数を計算する．


\section{凸最適化問題}

本稿では有限次元実Euclid空間 $\mathbb R^n$ を考え，
内積を $\langle\cdot,\cdot\rangle$ と書く．
凸集合の分離定理と微分可能な凸関数の基本事項を前提とする．

\begin{definition}
凸集合 $C\subset\mathbb R^n$ と凸関数
$f:C\to\mathbb R$ に対し，
\[
 \inf_{x\in C}f(x)
\]
を\emph{凸最適化問題}という．
最小値を達成する $x_*\in C$ を\emph{最小点}という．
\end{definition}

\begin{proposition}[微分可能な凸関数の一次条件]\label{prop:first-order}
$U\subset\mathbb R^n$ を開凸集合，
$f:U\to\mathbb R$ を微分可能とする．
$f$ が凸であることと
\[
 f(y)\ge f(x)+\langle\nabla f(x),y-x\rangle
 \qquad(x,y\in U)
\]
が成り立つことは同値である．
\end{proposition}
\begin{proof}
$f$ が凸なら $0<t<1$ について
\[
 f(x+t(y-x))\le(1-t)f(x)+tf(y)
\]
である．
両辺から $f(x)$ を引き $t$ で割って
$t\downarrow0$ とすれば
$\langle\nabla f(x),y-x\rangle\le f(y)-f(x)$ を得る．
逆に一次条件を $z=(1-t)x+ty$ に対して
$x$ と $y$ で適用すると
\[
 f(x)\ge f(z)+\langle\nabla f(z),x-z\rangle,\qquad
 f(y)\ge f(z)+\langle\nabla f(z),y-z\rangle.
\]
前者を $1-t$ 倍，後者を $t$ 倍して足すと，
$(1-t)(x-z)+t(y-z)=0$ だから
$(1-t)f(x)+tf(y)\ge f(z)$ となる．
従って $f$ は凸である．
\end{proof}

\begin{proposition}[最小点の存在]\label{prop:existence}
$C\subset\mathbb R^n$ を空でない閉集合，
$f:C\to\mathbb R$ を連続とする．
$f$ が $C$ 上で強制的，すなわち
$\norm x\to\infty$（$x\in C$）なら
$f(x)\to\infty$ を満たすとき，
$f$ は $C$ 上で最小値を達成する．
\end{proposition}
\begin{proof}
$x_0\in C$ を固定する．
$f$ は強制的だから
$S=\{x\in C:f(x)\le f(x_0)\}$ は有界である．
$C$ は閉で $f$ は連続なので $S$ は閉でもあり，
有限次元空間のHeine--Borelの定理から
$S$ はコンパクトである．
$S$ は $x_0$ を含むので空でない．
連続関数 $f|_S$ は最小値を達成し，
$C\setminus S$ の点では $f(x)>f(x_0)$ なので，
この点は $C$ 全体でも最小点である．
\end{proof}

\section{劣微分と法錐}

\begin{definition}
$U$ を開凸集合，$f:U\to\mathbb R$ を凸関数とする．
\[
 \partial f(x)=
 \{v\in\mathbb R^n:
 f(y)\ge f(x)+\langle v,y-x\rangle
 \text{ for all }y\in U\}
\]
を\emph{劣微分}という．
また $C\subset\mathbb R^n$ を凸集合，
$x\in C$ とするとき
\[
 N_C(x)=
 \{v\in\mathbb R^n:
 \langle v,y-x\rangle\le0
 \text{ for all }y\in C\}
\]
を\emph{法錐}という．
\end{definition}

\begin{proposition}[劣微分と勾配]\label{prop:grad-subdiff}
$f:U\to\mathbb R$ が凸で $x\in U$ で微分可能なら
\[
 \partial f(x)=\{\nabla f(x)\}.
\]
\end{proposition}
\begin{proof}
命題\ref{prop:first-order}から
$\nabla f(x)\in\partial f(x)$ である．
$v\in\partial f(x)$ とする．
任意の $h\in\mathbb R^n$ に対して
十分小さい正数 $t$ で $x\pm th\in U$ となる．
劣勾配不等式により
\[
 \frac{f(x+th)-f(x)}t\ge\langle v,h\rangle,
\]
\[
 \frac{f(x-th)-f(x)}{-t}\le\langle v,h\rangle
\]
となる．
$t\downarrow0$ で両辺の差商は
$\langle\nabla f(x),h\rangle$ に収束するので
$\langle v,h\rangle=\langle\nabla f(x),h\rangle$ である．
$h$ は任意だから $v=\nabla f(x)$ となる．
\end{proof}

\begin{theorem}[凸制約下の最適性条件]\label{thm:normal}
$C\subset\mathbb R^n$ を空でない凸集合とし，
$f:U\to\mathbb R$ を開凸集合 $U\supset C$ 上の
微分可能な凸関数とする．
$x_*\in C$ が $f$ の $C$ 上の大域的最小点であるための
必要十分条件は
\[
 -\nabla f(x_*)\in N_C(x_*)
\]
である．
\end{theorem}
\begin{proof}
$x_*$ が最小点とする．
任意の $y\in C$，$0<t<1$ に対して
$x_*+t(y-x_*)\in C$ だから
\[
 0\le f(x_*+t(y-x_*))-f(x_*).
\]
$t$ で割って $t\downarrow0$ とすると
$\langle\nabla f(x_*),y-x_*\rangle\ge0$ を得る．
従って $-\nabla f(x_*)\in N_C(x_*)$ である．
逆にこの法錐条件が成り立つなら，
命題\ref{prop:first-order}により
\[
 f(y)\ge f(x_*)+
 \langle\nabla f(x_*),y-x_*\rangle
 \ge f(x_*)
\]
が全ての $y\in C$ に対して成立する．
従って $x_*$ は大域的最小点である．
\end{proof}

\section{等式・不等式制約}

以下，$f,g_1,\dots,g_m:\mathbb R^n\to\mathbb R$ は
連続微分可能な凸関数，
$A:\mathbb R^n\to\mathbb R^r$ は線形写像，
$b\in\mathbb R^r$ とする．
主問題を
\[
 \begin{array}{ll}
 \text{最小化}&f(x),\\
 \text{制約}&g_i(x)\le0\quad(1\le i\le m),\\
             &Ax=b
 \end{array}
 \tag{P}
\]
と書く．

\begin{definition}
$\lambda=(\lambda_1,\dots,\lambda_m)\in[0,\infty)^m$，
$\nu\in\mathbb R^r$ に対し
\[
 L(x,\lambda,\nu)
 =f(x)+\sum_{i=1}^m\lambda_i g_i(x)
  +\langle\nu,Ax-b\rangle
\]
を\emph{Lagrange関数}という．
双対関数を
\[
 q(\lambda,\nu)=\inf_{x\in\mathbb R^n}
 L(x,\lambda,\nu)
\]
と定める．
\end{definition}

\begin{proposition}[弱双対性]\label{prop:weak-duality}
$x$ が (P) の実行可能点なら，
任意の $\lambda\ge0$，$\nu\in\mathbb R^r$ に対して
\[
 q(\lambda,\nu)\le f(x)
\]
が成り立つ．
\end{proposition}
\begin{proof}
実行可能点では $g_i(x)\le0$ かつ $Ax=b$ である．
従って
\[
 q(\lambda,\nu)
 \le L(x,\lambda,\nu)
 =f(x)+\sum_i\lambda_i g_i(x)
 \le f(x).
\]
これが主張である．
\end{proof}

\section{Slater条件と乗数の存在}

\begin{definition}
線形等式制約集合
\[
 S=\{x\in\mathbb R^n:Ax=b\}
\]
が空でなく，
ある $x_0\in S$ が存在して
\[
 g_i(x_0)<0\qquad(1\le i\le m)
\]
となるとき，\emph{Slater条件}が成立するという．
\end{definition}

\begin{lemma}[開凸集合による分離]\label{lem:strict-separation}
空でない開凸集合 $C\subset\mathbb R^d$ と
$x_0\notin C$ に対し，
零でない $v\in\mathbb R^d$ が存在して
\[
 \langle v,x\rangle\ge\langle v,x_0\rangle
 \qquad(x\in C)
\]
となる．
\end{lemma}
\begin{proof}
開凸集合とその外点の分離定理を適用すればよいが，
その帰結を明示する．
$D=C-x_0$ は原点を含まない開凸集合である．
有限次元のHahn--Banach分離定理によって
零でない線形汎関数 $\ell$ と実数 $\alpha$ が存在して
$\ell(y)\ge\alpha\ge0$（$y\in D$）となる．
$\ell(y)=\langle v,y\rangle$ と書き直せば
$\langle v,x-x_0\rangle\ge0$ が従う．
\end{proof}

\begin{theorem}[Slater条件とLagrange乗数]\label{thm:slater}
問題 (P) は最小点 $x_*$ を持ち，
その最小値 $p_*=f(x_*)$ は有限であるとする．
さらにSlater条件が成立すると仮定する．
このとき
$\lambda_i\ge0$ と $\nu\in\mathbb R^r$ が存在して，
\[
 \nabla f(x_*)+
 \sum_{i=1}^m\lambda_i\nabla g_i(x_*)+
 A^{\mathsf T}\nu=0,
 \qquad
 \lambda_i g_i(x_*)=0\quad(1\le i\le m)
\]
を満たす．
\end{theorem}
\begin{proof}
等式制約集合 $S=\{x:Ax=b\}$ を考える．
$\mathbb R^{m+1}$ の部分集合
\[
 C=\left\{(u_1,\dots,u_m,t):
 \begin{array}{l}
 \text{ある }x\in S\text{ が存在して}\\
 g_i(x)<u_i\quad(1\le i\le m),\\
 f(x)<t
 \end{array}\right\}
\]
を定める．
各点 $(u,t)\in C$ に対して不等式の余裕は正なので，
その点の十分小さい開近傍も $C$ に含まれる．
従って $C$ は開集合である．
また $f,g_i$ の凸性と $S$ のアフィン性によって，
二つの点を実現する $x,y\in S$ の凸結合を用いれば
$C$ は凸集合である．
$S$ は空でないので $C$ も空でない．
$(0,\dots,0,p_*)\notin C$ である．
実際，もし属すればある実行可能点 $x$ で
$f(x)<p_*$ が成り立ち，最小性に反する．

補題\ref{lem:strict-separation}から，
零でない $(\alpha_1,\dots,\alpha_m,\beta)$ が存在して
\[
 \sum_i\alpha_i u_i+\beta t\ge\beta p_*
 \qquad((u,t)\in C)
 \tag{1}
\]
となる．
$C$ は全ての座標について上方に閉じている：
$(u,t)\in C$ なら，任意の $h_i,h\ge0$ に対し
$(u_i+h_i,t+h)\in C$ である．
従って (1) において
$\alpha_i\ge0$，$\beta\ge0$ でなければならない．
Slater点 $x_0\in S$ に対して
$g_i(x_0)<0$ だから，
$u_i$ を負で $g_i(x_0)<u_i<0$ となるように取り，
$t>f(x_0)$ とすれば $(u,t)\in C$ である．
もし $\beta=0$ なら
零でない $\alpha\ge0$ について
$\sum_i\alpha_i u_i<0$ となるが，
(1) では $\sum_i\alpha_i u_i\ge0$ となり矛盾する．
従って $\beta>0$ である．

$\lambda_i=\alpha_i/\beta\ge0$ と置く．
任意の $x\in S$ に対して
$u_i\downarrow g_i(x)$，$t\downarrow f(x)$ として
(1) を極限に移せば
\[
 f(x)+\sum_i\lambda_i g_i(x)\ge p_*
 \qquad(x\in S)
 \tag{2}
\]
となる．
一方，$x_*$ は実行可能なので
$g_i(x_*)\le0$ であり，
\[
 f(x_*)+\sum_i\lambda_i g_i(x_*)\le p_*.
\]
(2) を $x=x_*$ に適用すると逆向きの不等式を得るので，
等号が成立する．
$\lambda_i g_i(x_*)\le0$ の総和が零であるから，
各項が零，すなわち
$\lambda_i g_i(x_*)=0$ となる．

さらに (2) の等号を $x_*$ が達成するので，
$x_*$ は
\[
 h(x)=f(x)+\sum_i\lambda_i g_i(x)
\]
のアフィン空間 $S$ 上の最小点である．
$S$ の方向空間は $\ker A$ なので，
任意の $z\in\ker A$ に対し
$t\mapsto h(x_*+tz)$ は $t=0$ で最小となる．
従って
$\langle\nabla h(x_*),z\rangle=0$
が全ての $z\in\ker A$ で成り立つ．
線形代数の基本恒等式
$(\ker A)^\perp=\operatorname{im}A^{\mathsf T}$
から，ある $\mu\in\mathbb R^r$ が存在して
$\nabla h(x_*)=A^{\mathsf T}\mu$ となる．
$\nu=-\mu$ と置けば
\[
 \nabla f(x_*)+\sum_i\lambda_i\nabla g_i(x_*)
 +A^{\mathsf T}\nu=0
\]
となり，結論を得る．
\end{proof}

\section{KKT条件の必要十分性}

\begin{definition}
実行可能点 $x_*$ と
$\lambda\in[0,\infty)^m$，$\nu\in\mathbb R^r$ が
次の条件を満たすとき，\emph{KKT条件}を満たすという：
\[
 \begin{aligned}
 &Ax_*=b,\qquad g_i(x_*)\le0,\\
 &\lambda_i\ge0,\qquad \lambda_i g_i(x_*)=0,\\
 &\nabla f(x_*)+\sum_i\lambda_i\nabla g_i(x_*)
   +A^{\mathsf T}\nu=0.
 \end{aligned}
\]
\end{definition}

\begin{theorem}[凸問題におけるKKT条件の十分性]
\label{thm:kkt-sufficient}
$f,g_i$ が微分可能凸関数なら，
KKT条件を満たす実行可能点 $x_*$ は
問題 (P) の大域的最小点である．
\end{theorem}
\begin{proof}
任意の実行可能点 $x$ を取る．
凸関数の一次条件から
\[
 f(x)\ge f(x_*)+
 \langle\nabla f(x_*),x-x_*\rangle
\]
である．
KKTの停留条件と $A(x-x_*)=0$ を用いると
\[
 \langle\nabla f(x_*),x-x_*\rangle
 =-\sum_i\lambda_i
 \langle\nabla g_i(x_*),x-x_*\rangle.
\]
また $g_i$ の凸性によって
\[
 \langle\nabla g_i(x_*),x-x_*\rangle
 \le g_i(x)-g_i(x_*).
\]
$\lambda_i\ge0$ と相補性
$\lambda_i g_i(x_*)=0$ から
\[
 \langle\nabla f(x_*),x-x_*\rangle
 \ge-\sum_i\lambda_i g_i(x)\ge0
\]
となる．従って $f(x)\ge f(x_*)$ であり，
$x_*$ は大域的最小点である．
\end{proof}

\begin{corollary}[Slater条件下のKKT判定]
\label{cor:kkt-iff}
Slater条件が成り立ち，問題 (P) の最小値が
ある実行可能点で達成されるとする．
このとき実行可能点 $x_*$ が最小点であることと，
ある $\lambda\ge0,\nu$ とともにKKT条件を満たすことは
同値である．
\end{corollary}
\begin{proof}
必要性は定理\ref{thm:slater}による．
十分性は定理\ref{thm:kkt-sufficient}による．
\end{proof}

\section{双対性と具体例}

\begin{theorem}[強双対性]\label{thm:strong-duality}
定理\ref{thm:slater}の仮定のもとで
\[
 \inf_{x\ {\rm feasible}}f(x)
 =\max_{\lambda\ge0,\ \nu\in\mathbb R^r}
 q(\lambda,\nu)
\]
が成り立つ．
\end{theorem}
\begin{proof}
定理\ref{thm:slater}によって
KKT乗数 $(\lambda_*,\nu_*)$ が存在する．
$L(x,\lambda_*,\nu_*)$ は $x$ について
微分可能な凸関数であり，
停留条件から
$\nabla_x L(x_*,\lambda_*,\nu_*)=0$
である．
命題\ref{prop:first-order}から
$x_*$ は $L(\cdot,\lambda_*,\nu_*)$ の
大域的最小点となる．
従って
\[
 q(\lambda_*,\nu_*)
 =L(x_*,\lambda_*,\nu_*)
 =f(x_*)+\sum_i\lambda_{*,i}g_i(x_*)
 =p_*,
\]
ここで最後の等号はKKTの相補性による．
弱双対性から全ての双対実行可能点で $q\le p_*$ なので，
$(\lambda_*,\nu_*)$ は双対問題の最大点となり，
両最適値が一致する．
\end{proof}

\begin{proposition}[等式制約とFenchel共役]\label{prop:fenchel}
$f:\mathbb R^n\to\mathbb R$ を凸関数，
$A:\mathbb R^n\to\mathbb R^r$ を線形写像とし，
等式制約 $Ax=b$ だけを考える．
この場合の双対関数は
\[
 q(\nu)=-f^*(-A^{\mathsf T}\nu)-\langle\nu,b\rangle
\]
となる．
\end{proposition}
\begin{proof}
定義により
\begin{align*}
 q(\nu)
 &=\inf_x\{f(x)+\langle\nu,Ax-b\rangle\}\\
 &=\inf_x\{f(x)+\langle A^{\mathsf T}\nu,x\rangle\}
   -\langle\nu,b\rangle\\
 &=-\sup_x\{\langle-A^{\mathsf T}\nu,x\rangle-f(x)\}
   -\langle\nu,b\rangle\\
 &=-f^*(-A^{\mathsf T}\nu)-\langle\nu,b\rangle.
\end{align*}
\end{proof}

\begin{example}[単一の不等式制約]
\[
 \text{最小化 }(x-2)^2\quad\text{制約 }x\le1
\]
を考える．最小点は $x_*=1$ であり，
KKT乗数は $\lambda=2$ である．
\end{example}
\begin{proof}
$f'(x)=2(x-2)$，$g(x)=x-1$，$g'(x)=1$ である．
$x_*=1$ では
$f'(1)+\lambda g'(1)=-2+\lambda=0$
となるため $\lambda=2$ である．
$\lambda\ge0$，$g(1)=0$ なのでKKT条件を満たす．
定理\ref{thm:kkt-sufficient}から最小点である．
また $x_0=0$ は $g(x_0)<0$ を満たすので
Slater条件も成立している．
\end{proof}

\begin{example}[Slater条件がない場合]
$f(x)=x$，$g(x)=x^2$ とし，
$g(x)\le0$ のもとで $f(x)$ を最小化する．
唯一の実行可能点は $x_*=0$ だが，
KKT乗数は存在しない．
\end{example}
\begin{proof}
$x^2\le0$ なら $x=0$ なので $x_*=0$ が唯一の実行可能点である．
停留条件は
\[
 f'(0)+\lambda g'(0)=1+\lambda\cdot0=0
\]
を要求するが，不可能である．
また $g(x)<0$ を満たす実数 $x$ は存在しないので
Slater条件は成立しない．
\end{proof}

\begin{remark}
KKT条件の必要性は制約の正則性に依存する．
本稿では有限次元の微分可能な凸問題に対し，
Slater条件から乗数の存在を分離定理によって証明した．
無限次元空間や非凸問題では，別の正則性条件と
より慎重な議論が必要である．
\end{remark}

\section{到達課題の完全解答}
\noindent\textbf{制約付き二次最適化のKKT計算．}
\begin{proposition}
\[
\min_{x\ge0}(x+1)^2
\]
の唯一の解は $x_*=0$ であり，
制約 $g(x)=-x\le0$ のKKT乗数は $\lambda_*=2$ である．
\end{proposition}
\begin{proof}
$f(x)=(x+1)^2$ は $x\ge0$ で導関数
$f'(x)=2(x+1)>0$ を持つから，
最小値は境界点 $x=0$ で一意に達成される．
Lagrange関数を $L(x,\lambda)=f(x)-\lambda x$
とすると，停留条件は
$2(x+1)-\lambda=0$，
双対実行可能性は $\lambda\ge0$，
相補性は $\lambda x=0$ である．
$(x,\lambda)=(0,2)$ はこれらを満たす．
凸問題で $x=1$ はSlater条件を満たすので，
KKT条件は最適性の十分条件である．
\end{proof}
\noindent\textbf{到達確認．}
$\min_{x\ge0}(x-1)^2$ の解とKKT乗数を求めよ．
\par\noindent\textbf{解答．}
最適解は内点 $x_*=1$，
相補性から $\lambda_*=0$ である．
\end{document}